\section{Compact spatial realization of the actual local third steering}
\label{trls:section}

The interval in Theorem~\ref{tr:local} has a complete compact
realization with the original periodic profile and physical diffusion.
We prove the substituted material, support, amplitude, and force
bounds here. Both older vorticities remain variables in this
realization. We also construct an explicit right continuation of the
same steering, which supplies a smooth compact time extension of its
forces. The required endpoint in this section is the local steering
endpoint, without an assertion about the full negative pulse.

\subsection{A proved interval, amplitude bounds, and material matrices}

Retain all the data of Theorem~\ref{tr:local}, and put
\begin{equation}\label{trls:times}
 \begin{gathered}
 \tau=\Gamma(t-t_1),\quad \tau_0=\sin s/8192,\quad
 t_*=t_1+\frac1{8192\sigma_2},\\
 \epsilon_R=\frac1{16384\sigma_2},\qquad
 I_*=[t_1,t_*],\quad I_+=[t_1,t_*+\epsilon_R].
 \end{gathered}
\end{equation}
Thus $0\leq\tau\leq3\tau_0/2$ on $I_+$, and the
original source steering remains its first ramp
$Z=Z_1(1-\eta(\tau))$ throughout. In particular every original
trial $m\in[0,3]$ gives the same solution on this interval.

The existence argument in Theorem~\ref{tr:local} extends to
this entire $I_+$. Indeed on the same closed cube of radius
$1/128$ about its actual six-dimensional entry state, the norm of
the scaled vector field is strictly less than $32/\sin s$.
On $0\leq\tau\leq3\tau_0/2$ its integral can therefore
move the state by at most
\[
 \frac{32}{\sin s}\frac{3\sin s}{2\cdot8192}
     =\frac3{512}<\frac1{128}.
\]
A first exit from that cube is impossible. Its fixed positive
denominator margins and bounded smooth derivatives extend the
unique solution to this closed interval, by the same integral-equation
argument. This proves a continuation of the actual source control,
including all its endpoint derivatives.

Write the dimensionless state as in \eqref{tr:scaled}, with
$E=E_1/\sigma_1^2$, $F_2^{\rm amp}=E_2/\sigma_2^2$,
$\omega=\Omega_1/\sigma_1$ and $V=\Omega_2/\sigma_2$.
To avoid confusing the second amplitude with the original periodic
profile $F$, denote $F_2^{\rm amp}$ by $F_a$ below. On all of
$I_+$ the already verified cube gives
\begin{equation}\label{trls:cube}
 \begin{gathered}
 63/64<h<4,\quad 11/16<g<6,\quad
 15/128<E<3,\quad 7/128<F_a<4,\\
 |\omega|<11,\quad |V|<1/16,\quad
 15/16<D=h^2+a^2<17,\\
 0\leq Z\leq Z_1<(9/8)\sin s<1/128,\\
 b=r_b\sin s<4\sin s,\quad 1\leq r_b\leq\sqrt{10}.
 \end{gathered}
\end{equation}
Use the exact $x,y,R,Y$ and phases in
\eqref{tr:coordinates}--\eqref{tr:inverse}. Then
\begin{equation}\label{trls:phasebounds}
 \frac1{64}<R=\frac{g^2+b^2}{D}<39,
 \qquad Y>1/9,\qquad 0<x<7,\quad 0<y<7.
\end{equation}
For the upper bound on $R$, use
$(36+16/512^2)/(15/16)<39$. Its lower bound and that of
$Y$ are the inequalities proved on the same cube in
Theorem~\ref{tr:local}. Positivity of $x,y$ is also proved
there; $x^2+y^2=R<39<49$ gives their upper bounds.
All three second components stay positive by the exact inverse,
without imposing a sign on either older vorticity.

We prove an explicit bound for the newest amplitudes. Set
\[
 X=-\Theta_3/P_3^*,\qquad
 W=-\sigma_2\Omega_3/(K_3P_3^*),\qquad
 \mathcal K=\frac{N_3}{\sigma_2^2\sin s\,R},\qquad
 \mathcal B=\frac Z{\sin s}.
\]
The exact pair \eqref{tr:physical_pair} becomes
\begin{equation}\label{trls:pair}
 X'=\mathcal K W-\delta_3RX,\qquad
 W'=\mathcal B X-\delta_3RW,
 \qquad \delta_3=dK_3^2/\Gamma.
\end{equation}
Its entry values satisfy $0<X(0)\leq e^4$ and
$0<W(0)\leq3e^4$ by the actual transition theorem.
The older geometry is already constructed on the compact interval
$I_+$, so the coefficient matrix of \eqref{trls:pair} is smooth
and bounded there. Its iterated integral series converges uniformly:
if its matrix norm is bounded by $B$, the term with $n$ integrals
is bounded by $(B\tau)^n/n!$ times the initial norm.
The series solves the integral equation, and the same iterated
estimate on the difference of two solutions proves uniqueness.
It therefore constructs the newest pair through all of $I_+$
independently of the following sign argument.
On the cube $N_3>0$, since $E_1,E_2,c,C,x,y,b$ are positive.
Removing the common damping and integrating the resulting positive
off-diagonal linear system proves $X,W>0$ throughout $I_+$.
More explicitly, a first zero of either positive coordinate would
contradict its positive initial value plus the integral of a
nonnegative multiple of the other coordinate.

There is a useful bound on $\mathcal K$ that retains the source
scale ratio. With $q=\sigma_1/\sigma_2=s/L$,
$\bar c=c/\sigma_1^2\leq1/4$ and
$\bar C=C/\sigma_1^2\leq a/32$, the exact numerator is
\[
 \frac{N_3}{\sigma_2^2\sin s}
 =\frac{q^2}{\sin s}\{(E+\bar c)y+\bar Cx\}
       +r_bF_a.
\]
The expression in braces is less than
$7(13/4+1/16384)<25$. Furthermore
\[
 \frac{q^2}{\sin s}=\frac{s^2}{L^2\sin s}
 \leq\frac{16s}{15L^2}
 \leq\frac{16}{15\cdot512\cdot16384^2},
 \qquad
 \frac{25\cdot16}{15\cdot512\cdot16384^2}<1.
\]
Here $\sin s\geq15s/16$ is the established source angle
estimate. Thus $0<N_3/(\sigma_2^2\sin s)<17$ and
\begin{equation}\label{trls:pairbox}
 0<\mathcal K<1088,\qquad 0\leq\mathcal B<9/8.
\end{equation}
Since $X,W>0$, addition of \eqref{trls:pair} gives
$(X+W)'\leq1089(X+W)$. Multiplication by
$e^{-1089\tau}$ and integration show
\[
 X+W\leq4e^4e^{1089\tau},\qquad
 1089\tau\leq\frac{3267\sin s}{16384}<\frac14.
\]
The power series gives $e^z\leq\sum_{n\geq0}z^n
=(1-z)^{-1}\leq4/3$ for $0\leq z\leq1/4$.
Consequently on the required interval and its proved right collar,
\begin{equation}\label{trls:amplitudes}
 0<-\Theta_3<6e^4P_3^*,\qquad
 |\Omega_3|<6e^4K_3P_3^*/\sigma_2.
\end{equation}
This uses the signed linear pair before any quotient estimate.

The exact matrices remain \eqref{tr:matrices}; all retain their
original birth times. With $h_0=\cos s_2$ and
$N=1+(3-h_0)/a$, define
\begin{equation}\label{trls:matrixbounds}
 \mathsf N_1=1,\qquad \mathsf N_2=2N,\qquad
 \mathsf N_3=32/\sin s.
 \qquad
 \|M_j\|,\|M_j^{-1}\|\leq\mathsf N_j.
\end{equation}
For the second matrix, the triangular factor and its inverse have
norm at most $1+|h-h_0|/a$. Since $h>0$, $h<4$, and
$h_0<1$, one has $|h-h_0|<4-h_0$. The difference
$2N-[1+(4-h_0)/a]=1+(2-h_0)/a$ is positive.
For the third, use $\mathcal L=(\zeta_2\ \zeta_3)$.
Its determinant is $-b$ and
$\|\mathcal L\|_F^2=D+R<56$, whereas
$\|\mathcal L(t_b)\|_F^2=r_b^2+1\leq11$.
The adjugate identity for two by two matrices therefore bounds
both products defining $M_3$ and its inverse by
$\sqrt{616}/b<32/\sin s$. This calculation allows both
older shears to evolve. The matrices have determinant one and
define the same global linear diffeomorphisms of $\mathbb R^2$.

\subsection{Support nesting with the unchanged source choice}

Let $C_M,C_N$ and $Y_3$ retain their exact values in
\eqref{tge:Y3} and \eqref{tts:sourceconstants}. Put
\[
 p=47Q_2/16-3,\qquad r=(31Q_2+1)/16.
\]
The literal outer $+1$ in \eqref{tge:Y3} gives the strict
stronger margins
\begin{equation}\label{trls:sourcemargins}
 \widehat\lambda_1^p\geq e^p(4C_NC_M),\qquad
 \widehat\lambda_1^r\geq e^r\sqrt{11}C_M/s_0.
\end{equation}
Indeed $\log\widehat\lambda_1\geq Y_3$ is at least one
plus either logarithmic summand divided by its positive exponent.
Also $N\leq C_N\widehat\lambda_1^{1/16}$ and
$\mathsf N_3\leq2C_M
\widehat\lambda_1^{(Q_2-1)/16}$, by
\eqref{tts:sourceconstants}. Multiplying these exact powers
proves
\begin{align*}
 \frac{2\mathsf N_2\mathsf N_3
                  \widehat\lambda_2^{-3}}
      {\widehat\lambda_1^{-3}}&\leq2e^{-p}<1,\\
 \frac{\sqrt{17}\mathsf N_3
                  \widehat\lambda_2^{-2}}{s_0}
     &\leq2\sqrt{17/11}\,e^{-r}<1.
\end{align*}
For the last strict inequalities, $Q_2\geq202$ implies
$p>1$ and $r>2$. The elementary inequality $e^z\geq1+z$
therefore gives $e^p>2$ and $e^r>3$, and
$2\sqrt{17/11}<3$ follows from $68<99$.
The earlier source margins give
\begin{equation}\label{trls:nestingnumbers}
 \begin{gathered}
 2\mathsf N_2\mathsf N_3\widehat\lambda_2^{-3}
      <\widehat\lambda_1^{-3},\quad
 \sqrt{17}\mathsf N_3\widehat\lambda_2^{-2}<s_0,\\
 \mathsf N_2\widehat\lambda_1^{-3}<s_0/(2C_\ell),
 \qquad \mathsf N_2\widehat\lambda_1^{-2}<s_0.
 \end{gathered}
\end{equation}
For the last two use $\mathsf N_2=2N$ and the proved
first-source margins $N\widehat\lambda_1^{-3}
<s_0/(4C_\ell)$ and $N\widehat\lambda_1^{-2}<s_0/2$.

Keep exactly the radii, periodic $F$, its even zero-mean primitive
$P$, and material cutoffs in \eqref{tts:envelopes}.
For every $t\in I_+$ a third support point satisfies
$|x|\leq\mathsf N_3\ell_3$. Thus
\[
 |M_2^{-1}x|\leq\mathsf N_2\mathsf N_3\ell_3<\ell_2/2,
 \qquad |K_2\zeta_2\cdot x|
    <K_2\sqrt{17}\mathsf N_3\ell_3<s_0.
\]
The entire second field is affine in a neighborhood of that point.
For every second support point the last two inequalities in
\eqref{trls:nestingnumbers} give $|x|<\ell_1/2$ and
$|K_1\zeta_1\cdot x|<s_0$. The first material matrix is
a rotation, so the first envelope is one there and its profile
is linear. Finally the original first support remains inside the
original affine cell of the radial base. The same inclusions hold
for every derivative support because their containing closed supports
are unchanged. All the inequalities are strict, so these are open
affine neighborhoods in the original physical coordinates.

\subsection{Full fields and exact scalar and vector equations}

Define the base and all three increments by
\eqref{tts:base} and \eqref{tts:fields}, using the exact
solution on $I_+$. All activation weights are one on this interval.
The fields agree at $t_1$, with every derivative, with the already
constructed fields through the third transition: the exact control
and phase jets match by the proof of Theorem~\ref{tr:local},
and repeated differentiation of the amplitude and material equations
matches their jets. They therefore give one smooth field on
$[0,t_*+\epsilon_R]$, with the original sine initial temperature
and zero initial velocity.

Use the exact scalar and full vector increments
\eqref{tts:Qrest}--\eqref{tts:spatialjets}, and the base
and total forces \eqref{tts:baseforce}. These formulas remain
the full residuals on $I_+$, as can be checked directly here.
On a $j$th support the proved older field is precisely
$G_{<j}\cdot x$ and $D_{<j}x$ with the unchanged sums in
\eqref{tts:material}. The identities
$\partial_t k_j=-D_{<j}^Tk_j$ and
$\partial_t M_j^{-1}=-M_j^{-1}D_{<j}$ give
\[
 (\partial_t+D_{<j}x\cdot\nabla)\xi_j=0,
 \qquad
 (\partial_t+D_{<j}x\cdot\nabla)g_j=0.
\]
In the scalar equation, the $F g_j$ part of
$v_j\cdot G_{<j}$ cancels the exact $\dot\Theta_j$
coupling, leaving the cutoff term, self-advection, modal damping,
and full Laplacian exactly as in \eqref{tts:scalarforce}.
In the vector equation, trace zero gives
$D_{<j}J+JD_{<j}^T=0$, whence the material derivative of
$J\nabla\Psi_j$ contributes $D_{<j}v_j$.
The reverse interaction $(v_j\cdot\nabla)(D_{<j}x)$
is a second $D_{<j}v_j$. Adding its other terms gives
\eqref{tts:vectorforce}, including both parent velocities,
the full buoyancy, and the moving-length term in $Q_{\rm r}$.
Adding the three identities counts every ordered older/newer
interaction once; outside their supports every derivative vanishes.
The identities therefore hold globally on $\mathbb R^2$.

Together with the base residual this proves exactly
\eqref{tts:PDE} on $[0,t_*+\epsilon_R]$, with $p=0$.
The streamfunction identity gives $\operatorname{div}u=0$.
Every spatial support remains in $B_{R_H}$ by the proved nesting.
The phase-coordinate Laplacian is still precisely
\eqref{tts:metric}, with its evolving positive matrix
$M_j^{-1}M_j^{-T}$ and all mixed derivatives. On a linear
profile cell the viscous force still includes
$-d\rho_jV_j$ and $-d\rho_jv_j$; the corresponding
Laplacians vanish there by \eqref{tts:core}. The full fields
retain the constant $P(0)$ in their cutoff region. Thus the
original profile has not been replaced by a modal eigenfunction.

\subsection{Actual all-order spatial and diffusion costs}

The following explicit data hold on the whole proved $I_+$:
\begin{equation}\label{trls:normdata}
 \begin{gathered}
 \mathsf r_{-,j}^2=(1,15/16,1/64)_j,\qquad
 \mathsf r_{+,j}^2=(1,17,39)_j,\quad
 \mathsf\kappa_j=K_j\mathsf r_{+,j},\\
 \mathsf U_{\Theta,j}=
 (3\sigma_1^2/K_1,\ 4\sigma_2^2/K_2,\ 6e^4P_3^*)_j,\\
 \mathsf U_{\Omega,j}=
 (11\sigma_1,\ \sigma_2/16,\ 6e^4K_3P_3^*/\sigma_2)_j,
 \qquad W_{j,1}=0,\\
 \mathsf A_*=693\sigma_1+12\Gamma
           +(81/8)\Gamma\sin s\,[\eta']_0.
 \end{gathered}
\end{equation}
The older data are the actual cube estimates, rather than their
smaller preceding-transition bounds. To verify
$|\dot\alpha|\leq\mathsf A_*$, use \eqref{tr:control}.
Here $Y>1/9$, $y<7$, $|\zeta_{1,1}|\leq1$,
$|\zeta_{2,1}|<5$, $b/D<(64/15)\sin s$,
$|\Omega_1|<11\sigma_1$, and
$|\Omega_2|<\sigma_2/16$. Therefore
\[
 |\mathcal F_3|
 <9\left(77\sigma_1+
   \frac{64\sin s}{15}\frac{\sigma_2}{16}\,5\right)
 =693\sigma_1+12\Gamma.
\]
The other control term has size
$|\dot Z|/Y\leq(81/8)\Gamma\sin s[\eta']_0$,
which proves the stated bound using the exact source derivative.

The following finite sums specify every spatial cost, with
$f_m,[F]_m,[P]_m,H_m,c_m$ defined as in Section~\ref{tts:section}:
\begin{align}\label{trls:costdefinitions}
 \mathsf D_j&=\mathsf A_*+\sum_{i<j}\mathsf U_{\Omega,i},&
 \mathsf G_j&=A_0+\sum_{i<j}K_i\mathsf U_{\Theta,i}
                                      \mathsf r_{+,i},\nonumber\\
 \mathsf Q_j&=\frac{\mathsf U_{\Omega,j}}
                         {K_j^2\mathsf r_{-,j}^2},&
 \mathsf Q_{{\rm r},j}&=
 \frac{\mathsf\kappa_j\mathsf U_{\Theta,j}
                   +2\mathsf D_j\mathsf U_{\Omega,j}}
      {K_j^2\mathsf r_{-,j}^2},\\
 \mathsf e_{j,m}&=(\mathsf N_j/\ell_j)^mf_m,&
 \mathsf B_{j,m}(W)&=\sum_{q=0}^m\binom mq
          \mathsf\kappa_j^q[W]_q\mathsf e_{j,m-q},\nonumber\\
 \mathsf C_{j,m}(P)&=\sum_{q=0}^m\binom mq
          \mathsf\kappa_j^q[P]_q\mathsf e_{j,m-q+1}.&&\nonumber
\end{align}
In the occurrences $\mathsf B_{j,m}(W)$ the argument is a
periodic test profile, not the signed amplitude $W$ of
\eqref{trls:pair}. The operator norm of
$J\zeta_i\otimes\zeta_i/r_i^2$ equals one, so
$\|D_{<j}\|\leq\mathsf D_j$ with the entire second
vorticity included. The gradient sum similarly gives
$|G_{<j}|\leq\mathsf G_j$. Differentiating
$\rho_j$ gives $|\dot\rho_j/\rho_j|\leq2\mathsf D_j$.
These prove the bounds $|Q_j|\leq\mathsf Q_j$ and
$|Q_{{\rm r},j}|\leq\mathsf Q_{{\rm r},j}$.
The product rule with its $\binom mq$ allocations and the
linear material bound \eqref{trls:matrixbounds} proves
$|W(\xi_j)g_j|_m\leq\mathsf B_{j,m}(W)$ and
$|P(\xi_j)\nabla g_j|_m\leq\mathsf C_{j,m}(P)$.

In particular the complete damping and diffusion terms obey
\begin{align}\label{trls:diffusioncost}
 \mathsf D_{\theta,j,m}
 &=d\mathsf U_{\Theta,j}
   \{\mathsf\kappa_j^2\mathsf B_{j,m}(F)
                     +2\mathsf B_{j,m+2}(F)\},\nonumber\\
 \mathsf D_{u,j,m}
 &=d\mathsf Q_j
   \{\mathsf\kappa_j^2\mathsf B_{j,m+1}(P)
                     +2\mathsf B_{j,m+3}(P)\}.
\end{align}
The two coordinate traces of the Laplacian account for the factor
two in each expression. Direct termwise application to the full
residuals proves the bounds
\begin{align}\label{trls:forcecost}
 \mathsf E_{\theta,j,m}={}&
 \mathsf Q_j\mathsf G_j\mathsf C_{j,m}(P)
                     +\mathsf D_{\theta,j,m}\nonumber\\
 &+\mathsf Q_j\mathsf U_{\Theta,j}
    \sum_{q=0}^m\binom mq
       \mathsf B_{j,q+1}(P)\mathsf B_{j,m-q+1}(F),\nonumber\\
 \mathsf E_{u,j,m}={}&
 (2\mathsf D_j\mathsf Q_j+\mathsf Q_{{\rm r},j})
                  \mathsf B_{j,m+1}(P)
 +\mathsf U_{\Theta,j}\mathsf B_{j,m}(F)
 +\mathsf D_{u,j,m}\nonumber\\
 &+\mathsf Q_j^2\sum_{q=0}^m\binom mq
       \mathsf B_{j,q+1}(P)\mathsf B_{j,m-q+2}(P),\\
 |E_{\theta,j}|_m\leq{}&\mathsf E_{\theta,j,m},\qquad
 |E_{u,j}|_m\leq\mathsf E_{u,j,m}.\nonumber
\end{align}
Indeed the scalar product sum bounds
$v_j\cdot\nabla V_j$; the vector sum bounds
$(\nabla v_j)v_j$, allocating all $m$ derivatives in
each case. The remaining terms are exactly the linear terms of
\eqref{tts:scalarforce} and \eqref{tts:vectorforce}.
There is no activation derivative on $I_+$ because all weights
are one.

Let $B_{0,m}(F)=\sum_{q=0}^m\binom mq\mu^q[F]_qc_{m-q}$.
The control constant $\mathsf A_2$ bounding $|\ddot\alpha|$
is explicitly defined below. The base and sums now give global norms
on the original $\mathbb R^2$:
\begin{align}\label{trls:globalnorms}
 |\theta|_m&\leq(A_0/\mu)B_{0,m}(F)
       +\sum_{j=1}^3\mathsf U_{\Theta,j}\mathsf B_{j,m}(F),
 \nonumber\\
 |u|_m&\leq\mathsf A_*H_m
       +\sum_{j=1}^3\mathsf Q_j\mathsf B_{j,m+1}(P),\nonumber\\
 |f_\theta|_m&\leq2d(A_0/\mu)B_{0,m+2}(F)
                       +\sum_{j=1}^3\mathsf E_{\theta,j,m},\\
 |f_u|_m&\leq\mathsf A_2H_m
       +\mathsf A_*^2\sum_{q=0}^m\binom mqH_{q+1}H_{m-q}
       +(A_0/\mu)B_{0,m}(F)+2d\mathsf A_*H_{m+2}
                       +\sum_{j=1}^3\mathsf E_{u,j,m}.\nonumber
\end{align}
These follow from the displayed fields and base residual; the
initial interpolation is already complete on $I_+$. They bound
the entire fields and forces, not only their common affine cell.
For each fixed bound in this display, integration over the required
interval multiplies it by $1/(8192\sigma_2)$; over $I_+$ the
factor is $3/(16384\sigma_2)$. The earlier proved norms together
with these, taking the maximum on their two adjacent intervals,
give the corresponding bounds through $t_*$.

The exact retained localization and newest diffusion factors include
\begin{equation}\label{trls:frequencycost}
 \begin{gathered}
 \mathsf N_1/\ell_1=C_\ell\mu/s_0,\quad
 \mathsf N_2/\ell_2=2N\mu\widehat\lambda_1^3,\quad
 \mathsf N_3/\ell_3=
                32\mu\widehat\lambda_2^3/\sin s,\\
 d\mathsf U_{\Theta,3}K_3^2
       =6e^4dA_0\mu\widehat\lambda_3^{9/8},\quad
 d\mathsf Q_3K_3^3
       =384e^4dA_0\mu\widehat\lambda_3^{9/8}/\sigma_2.
 \end{gathered}
\end{equation}
They are explicit factors in upper force costs. The exact three
modal exponents over $I_*$ are
\begin{equation}\label{trls:exponents}
 dK_1^2|I_*|,\qquad dK_2^2\int_{I_*}D\,dt,\qquad
 dK_3^2\int_{I_*}R\,dt.
\end{equation}
The second is between $(15/16)dK_2^2|I_*|$ and
$17dK_2^2|I_*|$, and the third between
$dK_3^2|I_*|/64$ and $39dK_3^2|I_*|$.
The exact newest gain is \eqref{tr:costs} with
$\tau_f=\tau_0$; all divisions there are valid by positivity
of the pair just proved. Both older numerator-loss terms in
\eqref{tr:losses} remain present, separately from these modal
exponents.

\subsection{Exact source jets, mixed costs, and compact time forces}

We give an explicit finite recursion for all control derivatives.
Let $z=(h,g,E,F_a,\omega,V,X,W)$ and introduce the switch jets
$\lambda^{[j]}$. Define $Z=Z_1(1-\lambda^{[0]})$,
and use the exact inverse \eqref{tr:inverse} with this $Z$.
The eight-dimensional vector field $\mathcal V^{\rm loc}$
uses the coordinate right-hand sides in \eqref{tr:scaled}
with $F=F_a$, ordered as $(h,g,E,F_a,\omega,V)$,
followed by the two right-hand sides in \eqref{trls:pair}.
In particular it depends
on $\lambda^{[0]}$, not its derivatives. Define the exact
physical angular function
\begin{equation}\label{trls:controlfunction}
 \Phi(z,\lambda^{[0]},\lambda^{[1]})
  =\frac{\sigma_1\omega y\zeta_{1,1}
           +(b\sigma_2V/D)\zeta_{2,1}
           +\Gamma Z_1\lambda^{[1]}}{Y}.
\end{equation}
Substituting $\lambda^{[j]}=\eta^{(j)}(\tau)$ gives
exactly $\dot\alpha$ in \eqref{tr:control}, because
$\dot Z=-\Gamma Z_1\eta'(\tau)$.
Thus every source-ramp derivative retains its original factor
$\Gamma^rZ_1$. No constant switch is used on the collar.

On the explicit compact rectangle
\begin{equation}\label{trls:jetbox}
 \begin{gathered}
 h\in[63/64,4],\quad g\in[11/16,6],\quad
 E\in[15/128,3],\quad F_a\in[7/128,4],\\
 \omega\in[-11,11],\quad V\in[-1/16,1/16],\quad
 X,W\in[0,6e^4],\quad \lambda^{[0]}\in[0,1],
 \end{gathered}
\end{equation}
the estimates for $D,R,Y$ used above are uniform and positive.
The same elementary inequalities apply at its boundary, since their
final margins are strict even with all displayed weak bounds.
The vector field and $\Phi$ are consequently smooth on a
neighborhood of this rectangle for any finite control-jet bounds.
Define the derivation
\begin{equation}\label{trls:jetderivation}
 \mathcal C_{\rm loc}
 =\sum_{l=1}^8\mathcal V_l^{\rm loc}
                   \partial_{z_l}
   +\sum_{j\geq0}\lambda^{[j+1]}
                   \partial_{\lambda^{[j]}}.
\end{equation}
Each application to a finite expression uses only finitely many
terms. The chain rule gives, by induction on $n$,
\begin{equation}\label{trls:controljets}
 \begin{gathered}
 \alpha^{(n+1)}=
       \Gamma^n\mathcal C_{\rm loc}^n\Phi
       \big|_{\lambda^{[j]}=\eta^{(j)}(\tau)},\\
 \mathsf A_{n+1}=\Gamma^n\max
 \left\{|\mathcal C_{\rm loc}^n\Phi|:
    z,\lambda^{[0]}\text{ in \eqref{trls:jetbox}},
    \ |\lambda^{[j]}|\leq[\eta^{(j)}]_0,
                            \ 1\leq j\leq n+1\right\},\quad n\geq0.
 \end{gathered}
\end{equation}
These are finite maxima of specified elementary expressions on
specified compact rectangles. They contain no supremum over a
future trajectory. The preceding identity proves
$|\alpha^{(n+1)}|\leq\mathsf A_{n+1}$; the sharper
$\mathsf A_*$ in \eqref{trls:normdata} may be used at first
order. In particular the base constant in
\eqref{trls:globalnorms} is precisely $\mathsf A_2$.

For clarity, the induction behind the displayed identity differentiates
a function of $z$ along its eight exact state rows and differentiates
each $\eta^{(j)}$ to $\eta^{(j+1)}$. Physical time
adds a factor $\Gamma$ at each step. Conversely substituting
the state scalings into these rows recovers the six older physical
equations and newest physical pair; their inverse scalings are
exact and all scale denominators are positive.

All mixed space and time norms are then evaluated by the complete
pullback operators \eqref{tts:pullbackops} and their identity
\eqref{tts:pullbackidentity}, replacing the transition derivation
by \eqref{trls:jetderivation} and its rectangle by
\eqref{trls:jetbox}. Explicitly start from the zeroth constants
\eqref{trls:matrixbounds} and \eqref{trls:normdata} and
differentiate
\[
 \dot M_j=D_{<j}M_j,\qquad
 \dot A_j=-A_jD_{<j},\qquad
 \dot k_j=-D_{<j}^Tk_j,\qquad A_j=M_j^{-1}/\ell_j.
\]
For example
$\partial_t^{n+1}A_j=-\sum_{q=0}^n\binom nq
(\partial_t^qA_j)(\partial_t^{n-q}D_{<j})$.
The product rule gives the other two recursions with the displayed
matrix order and signs. Coefficient jets of the parent sums are
given by \eqref{trls:jetderivation}; all activation jets on
$I_+$ vanish. On every layer support bound each factor $x^\beta$
in the pullback expression by
$(\mathsf N_j\ell_j)^{|\beta|}$, each periodic derivative
by its original $[F]_r$ or $[P]_r$, and each envelope derivative
by $f_r$. Sum the resulting finite absolute products with the
binomial coefficients produced by these recursions. The chain and
product rules prove the resulting bound for every prescribed mixed
order. On the base use $R_H$, its fixed profile derivatives, and
\eqref{trls:controljets} instead. The recursion has finite depth
equal to the requested derivative order and retains every time
derivative of the actual Laplacian metric and modal diffusion.

Finally extend the original sine datum constantly to negative time
with zero velocity as in Section~\ref{tts:section}, and take
$\epsilon_L=1/\nu$. Let $f_\bullet^{\rm raw}$ be the full
residuals \eqref{tts:rawcollar} of these fields on
$[-\epsilon_L,t_*+\epsilon_R]$. On the nonnegative interval
they equal the proved telescoping force formulas. Define
\begin{equation}\label{trls:timecutoff}
 \beta_*(t)=\eta(1+t/\epsilon_L)
             \eta(1+(t_*-t)/\epsilon_R),\qquad
 \widetilde f_\bullet=\beta_*f_\bullet^{\rm raw}.
\end{equation}
Here the right collar is the actual source steering continuation
already constructed, with physical length $\epsilon_R$ in
\eqref{trls:times}. The cutoff equals one on $[0,t_*]$;
flatness of the original step makes extension by zero smooth at
both outer times. The forces are therefore smooth and compact in
space and time, and solve the exact PDE on the entire required
$[0,t_*]$. Their complete cutoff costs are
\begin{equation}\label{trls:cutoffcost}
 \begin{gathered}
 \mathsf B_{\beta,r}=\sum_{q=0}^r\binom rq
      \epsilon_L^{-q}\epsilon_R^{-(r-q)}
       [\eta^{(q)}]_0[\eta^{(r-q)}]_0,\\
 \|\partial_t^n\partial_x^\gamma\widetilde f_\bullet\|_\infty
 \leq\sum_{q=0}^n\binom nq\mathsf B_{\beta,q}
  \|\partial_t^{n-q}\partial_x^\gamma
                                f_\bullet^{\rm raw}\|_\infty.
 \end{gathered}
\end{equation}
The raw mixed norms on $I_+$ have the explicit recursion just
proved, those on the preceding slab have the proved earlier
recursions, and those on the constant left extension are obtained
by differentiating its fixed sine datum and residual. Taking their
maximum gives every norm in \eqref{trls:cutoffcost}.

Oddness of the original scalar profile, evenness of its primitive
and cutoffs, and the linear material maps give the same odd
temperature, velocity and forces and even vorticity as before.
The compact smooth streamfunction has finite-energy velocity and
still satisfies the exact Biot--Savart identity proved in
Section~\ref{tts:section}. Thus the local source steering, its
moving endpoint, and its proved right collar have a complete
compact original-profile PDE realization with all physical
diffusion, localization, and control costs specified.
